\section{Ledger Architecture}
\label{sec:ledger}

% Core formal contribution #4: the minimal verification ledger.

% §6.1 Design principles
%   The ledger is motivated by verification requirements, not by
%   analogy to financial blockchains. It records two kinds of entries:
%   (a) claim verification records (commitment + result), and
%   (b) governance decisions (benchmark proposals + votes + outcomes).
%
%   Required properties:
%   - Append-only (entries cannot be deleted or modified)
%   - Publicly auditable (any participant can verify the chain)
%   - Substrate-distributed (replicated across substrates; no single
%     substrate controls the ledger)
%
%   NOT required:
%   - Token economics or financial incentives
%   - Proof-of-work or proof-of-stake consensus
%   - Smart contract execution
%   - High transaction throughput (governance decisions are infrequent;
%     claim verifications scale with compute, not consensus)
%
%   The minimal construction is an authenticated append-only data
%   structure (Merkle tree or similar) replicated across substrate-
%   exclusive nodes.

% §6.2 Verification entries
%   Each verification entry contains:
%   - Agent identifier a_j
%   - Commitment Com(C, r)
%   - Algorithmic witness identifier W⊥_alg
%   - Benchmark/protocol identifier B (must be governance-accepted)
%   - Verification result (pass/fail + timestamp + raw data hash)
%   - Opened commitment (C, r) for public audit
%   Signed by both the agent (commitment phase) and the witness
%   (verification phase).

% §6.3 Governance entries
%   Each governance entry contains:
%   - Benchmark proposal (specification, target capability, proposer)
%   - Deliberation record (analysis submitted during the window)
%   - Vote record (per-validator votes, substrate distribution)
%   - Outcome (accepted/rejected + effective date for accepted forks)
%   The governance record is permanent: rejected proposals remain in
%   the ledger with their deliberation records. This provides an
%   audit trail for benchmark evolution and prevents re-proposal of
%   previously-studied-and-rejected benchmarks without new evidence.

% §6.4 Consensus among substrate-exclusive nodes
%   The ledger is replicated across nodes on different substrates.
%   For verification entries: append requires agreement from the
%   committing agent and the verifying witness (two-party, cross-
%   substrate). No general consensus needed for per-claim entries.
%   For governance entries: acceptance requires a threshold of
%   cross-substrate validators (e.g., supermajority with at least
%   one validator per substrate). This is a special-purpose consensus
%   that is simpler than general Byzantine consensus because the
%   entries are governance decisions (infrequent, high deliberation)
%   rather than financial transactions (frequent, low latency).

% §6.5 Relation to the GFM harness ledger
%   The GFM harness already maintains a ledger (session/ledger.json)
%   that tracks agents, capabilities, trust, and cooperative
%   capabilities. Paper 5's construction provides the cryptographic
%   backbone that makes this ledger tamper-evident, cross-substrate
%   verifiable, and governance-auditable.

% TODO: formal definitions, data structures, consensus specification.
